
\begin{exercise}
In this problem, boxes can be empty. How many ways are there to distribute 6 balls in to 3 boxes if\\
a. both the balls and boxes are labeled?\\
b. balls are labeled, but the boxes are unlabeled?\\
c. the balls are unlabeled, but the boxes are labeled?\\
d. both the boxes and balls are unlabeled?
\end{exercise}

\begin{exercise}
In this problem, boxes must be non-empty. How many ways are there to distribute 8 balls in to 4 boxes if\\
a. both the balls and boxes are labeled?\\
b. balls are labeled, but the boxes are unlabeled?\\
c. the balls are unlabeled, but the boxes are labeled?\\
d. both the boxes and balls are unlabeled?
\end{exercise}

\begin{exercise}
How many 8-bit strings begin with a 010 or begin with 11?  How many 8-bit strings begin with a 010 or end with 11?
\end{exercise}

\begin{exercise}
Find the number of solutions to $x+y+z=100$ in non-negative integer $x$, $y$, and $z$, where $x \le 50$, $y \le 40$, and $z \le 30$.
\end{exercise}

\begin{exercise}
How many permutations of the numbers 1 through 8 are there in which none of the terms 12, 34, 56, 78 appear?
\end{exercise}

\begin{exercise}
How many permutations of the numbers 1 through 8 are there in which none of the even integers remain in its natural position?
\end{exercise}

\begin{exercise}
Do the enigma machine example from class where the number of rotors is 8, rather than 5.
\end{exercise}

\begin{exercise}
How many permutations of 1,1,2,2,3,3,4,4,5,5 are there in which no two adjacent numbers are equal?
\end{exercise}

\begin{exercise}
Suppose you roll a 6 sided die 12 times. In how many ways can all 6 six numbers appear?
\end{exercise}

\begin{exercise}
How many integers in the set of integers from 1 to 1000 are divisible by a prime number less than 10?
\end{exercise}

\begin{exercise} 
A bridge hand consists of 13 cards from a standard 52 card deck.  How many different hands contain at least on card from each of the four suits?
\end{exercise}

\begin{exercise} 
Let $A=\{1,2,3,4,5\}$. How many injective functions $f:A\rightarrow A$  have the property that for each $x$ in $A$, $f(x)\ne x$ ?
\end{exercise}

\begin{exercise} 
Suppose you planned on giving 7 gold stars to some of the 13 star students in your class. Each student can receive at most one star. How many ways can you do this?
\end{exercise}

\begin{exercise} 
What is the formula for the number of ways to place $n$ unlabeled balls into $k$ unlabeled boxes {\bf without} the restriction that no box may be empty?
\end{exercise}

\begin{exercise} 
What is the formula for the number of ways to place $n$ labeled balls into $k$ unlabeled boxes {\bf without} the restriction that no box may be empty?
\end{exercise}


%\begin{proofp}
%Consider two sets $X$ and $Y$ with $|X|=m$ and $|Y|=n$ A function $f:X \rightarrow Y$ is called {\bf almost-onto} if $f$ misses at most one element of $Y$. Find and prove a formula for the number of almost onto functions from $X$ to $Y$.
%\end{proofp}

\begin{proofp}
Prove by induction $S(n,3)>3^{n-2}$ for all $n\ge 6$.
\end{proofp}


\begin{proofp} $\bigstar$Prove
$$(1-a_1)(1-a_2) \cdots (1-a_k)= 1- (a_1+a_2+ \cdots+a_k)+ (a_1a_2+\cdots +a_{k-1}a_{k})+ \cdots +(-1)^k(a_1 \cdots a_k) $$
\end{proofp}

\begin{proofp}
For fixed $m$, find and prove the number of Hamiltonian Paths on on $K_{m,n}$ as a function of $n$.
\end{proofp}

\begin{proofp} $\bigstar$
From the text, Exercise 4.6
\end{proofp}

\begin{proofp}$\bigstar$
From the text, Exercise 4.7 
\end{proofp}


\begin{proofp}$\bigstar$
Let S(n,k) be the Stirling numbers of the second kind. Prove \\
a. $S(n, n-1)= \binom{n}{2}$\\
b. $S(n, n-2)= \binom{n}{3}+3\binom{n}{4}$
\end{proofp}



\begin{proofp}$\bigstar$
Let S(n,k) be the Stirling numbers of the second kind. Prove $S(n,k)=\sum_{m=k-1}^{n-1} \binom{n-1}{m} S(m,k-1)$
\end{proofp}



\begin{proofp}$\bigstar$
Show that if $n$ is a positive integer, then
$$ n!=\binom{n}{0}d_n+\binom{n}{1}d_{n-1}+\binom{n}{2}d_{n-2}+ \cdots +\binom{n}{n-1}d_1 +\binom{n}{n}d_0$$
where $d_n=(n-1)(d_{n-1}+d_{n-2})$ with $d_1=0$ and $d_2=1$ is the number of derangements of $n$ objects.
\end{proofp}

\begin{proofp}$\bigstar$
Let $A$ be a subset of $B$ with $|A|=m$ and $|B|=n$, and let $r$ be an integer $m \le r \le n$.  Show the number of $r$ element subsets of $B$ which contain $A$ as a subset is $\binom{n-m}{r-m}$.
\end{proofp}


\begin{proofp}$\bigstar$
Use the Inclusion-Exclusion Principle to show that for any non-negative integers $m,r,n$ such that $m \le r \le n$,

$$\binom{n-m}{r-m}=\sum^{m}_{i=0}(-1)^i \binom{m}{i}\binom{n-i}{r}.$$
\end{proofp}